\section{TDA2030 y Criterios de Diseño}
\subsection{Justificación de la Selección del TDA2030}
El TDA2030 es un circuito integrado monolítico encapsulado en TO-220, concebido para su uso como amplificador de baja frecuencia en clase AB. Es ampliamente utilizado en entornos educativos e industriales debido a:
\begin{itemize}[label=\small\color{primary}$\blacksquare$]
    \item Capacidad de entregar hasta 14 W de potencia de salida sobre una carga de $4\,\Omega$.
    \item Amplio rango de tensión de alimentación (hasta $\pm 18$\,V en fuente simétrica o 36\,V en fuente simple).
    \item Protecciones integradas contra cortocircuitos y sobretemperatura.
    \item Bajo costo y requerimiento de pocos componentes externos.
\end{itemize}

\subsection{Análisis del Circuito Estándar (Datasheet)}
El diseño adoptado corresponde al esquema típico provisto por el fabricante con alimentación simple (o simétrica, según corresponda). 

Las ecuaciones de la ganancia en tensión en banda media vienen dadas por la red de realimentación negativa:
\begin{equation}
    A_v = 1 + \frac{R_3}{R_2}
\end{equation}

Asimismo, se consideran las frecuencias de corte de los filtros de desacoplo de entrada y salida:
\begin{equation}
    f_{c,in} = \frac{1}{2\pi \cdot R_{in} \cdot C_{in}}
\end{equation}